Après avoir choisi un ordre pour les inconnues scalaires,
\[
\mathbf u=
\bigl((u_K)_{K\in\mathfrak M},
      (u_{K^\ast})_{K^\ast\in\mathfrak M^\ast}\bigr)^\top
\in\mathbb R^{\lvert\mathfrak M\rvert+\lvert\mathfrak M^\ast\rvert},
\]
et un ordre pour les diamants, le gradient discret s'écrit
\[
\mathbf g=\gradmat_{\ensdiamantplat}\mathbf u,
\qquad
\mathbf g=
\bigl((\grad^D u)_x,(\grad^D u)_y\bigr)_{D\in\ensdiamantplat}^{\top}.
\]

Pour un diamant \(D=D_{\sigma,\sigma^\ast}\), avec
\[
\sigma=[x_K,x_L],
\qquad
\sigma^\ast=[x_{K^\ast},x_{L^\ast}],
\]
posons
\[
\mathbf a_D
=
\frac{\vect{\nu}_{\sigma K}}
{\sin(\angle D)\,\mesure(\sigma^\ast)},
\qquad
\mathbf b_D
=
\frac{\vect{\nu}_{\sigma^\ast K^\ast}}
{\sin(\angle D)\,\mesure(\sigma)}.
\]
La formule du gradient devient alors
\[
\grad^D u
=
\mathbf a_D(u_L-u_K)
+
\mathbf b_D(u_{L^\ast}-u_{K^\ast}).
\]

Le bloc de deux lignes de \(\gradmat_{\ensdiamantplat}\) associé à \(D\) ne comporte donc que quatre colonnes non nulles :
\[
\boxed{
\left(\gradmat_{\ensdiamantplat}\right)_{D,\{K,L,K^\ast,L^\ast\}}
=
\begin{pmatrix}
-\mathbf a_D & \mathbf a_D & -\mathbf b_D & \mathbf b_D
\end{pmatrix}}
\]
où chaque entrée est un vecteur colonne de \(\mathbb R^2\).

En coordonnées cartésiennes, si
\[
\mathbf a_D=
\begin{pmatrix}a_{D,x}\\a_{D,y}\end{pmatrix},
\qquad
\mathbf b_D=
\begin{pmatrix}b_{D,x}\\b_{D,y}\end{pmatrix},
\]
ce bloc est explicitement
\[
\begin{pmatrix}
-a_{D,x} & a_{D,x} & -b_{D,x} & b_{D,x}\\
-a_{D,y} & a_{D,y} & -b_{D,y} & b_{D,y}
\end{pmatrix}.
\]

La matrice globale est obtenue en empilant ces blocs pour tous les diamants :
\[
\gradmat_{\ensdiamantplat}
=
\begin{pmatrix}
\gradmat_{D_1}\\
\vdots\\
\gradmat_{D_{\lvert\ensdiamantplat\rvert}}
\end{pmatrix}
\in
\mathbb R^{2\lvert\ensdiamantplat\rvert
\times
(\lvert\mathfrak M\rvert+\lvert\mathfrak M^\ast\rvert)}.
\]
Chaque bloc contient au plus huit coefficients scalaires non nuls, ce qui fait de \(\gradmat_{\ensdiamantplat}\) une matrice creuse.
